This dissertation investigates the mechanisms of surplus distribution between formal and informal sectors in India's textile and apparel value chains. It addresses the debate between the "displacement" view of modernization and the "adverse incorporation" view of structural exploitation. Using a state-industry panel dataset for 2010--2024, I integrate formal sector data from the Annual Survey of Industries (ASI) with informal sector data from the Annual Survey of Unorganized Sector Enterprises (ASUSE).

The empirical results provide robust evidence for an "adverse incorporation" equilibrium. I find that formal productivity growth is positively associated with informal outsourcing ($\beta = 9.07 \times 10^{-4}, p=0.043$), an effect that is 4.4 times larger in high-enforcement states, indicating strategic regulatory arbitrage. Crucially, there is a statistical zero wage pass-through to informal workers ($\beta = -0.014, p=0.636$), confirming that formal lead firms capture the entirety of productivity rents. Furthermore, informal employment exhibits extreme structural persistence ($\beta = 0.8503, p<0.001$). Empirically-grounded simulations demonstrate that conventional labor enforcement perversely scales the informal sector without improving earnings, while only direct wage floors successfully break the monopsony trap. These findings demonstrate that formal modernization suppresses rather than dissolves informal activity, necessitated by the structural power imbalances in globalized production networks.
